Fix $\tau$ and suppress it. Along $\gamma_{y}$ write the Hubble-normalised generator as $G_{y}(s) = H_{y}(s) + \mathcal{E}_{y}(s)$, with the model part $H_{y}(s) = -\sum_{i}u^{i}_{y}(s)\,p_{i}B_{i}$ and the error $\mathcal{E}_{y}(s) = E_{\tau}(\gamma_{y}(s))[u_{y}(s)]$, as in the admissibility display of Definition \claimref{obj.interior.weighted-blow-up-hubble-normalisation-admissible}. Then $T_{y}' = G_{y}T_{y}$, $(T^{0}_{y})' = H_{y}T^{0}_{y}$, both equal to $1$ at $s = 0$. The bounds used are those of the proof of Theorem \claimref{interior.transport-level-convergence-along-tangential-curves}: $\|H_{y}\| \le 1$ (Cauchy-Schwarz with $\sum_{i}p_{i}^{2} = 1$ and $\|B_{i}\| = 1$), $\|\mathcal{E}_{y}\| \le \varepsilon$, $\|G_{y}\| \le 1 + \varepsilon$, hence $\|T_{y}(s)\| \le e^{s(1+\varepsilon)}$, $\|T^{0}_{y}(s)\| \le e^{s}$, $\|T_{y}(\lambda)T_{y}(s)^{-1}\| \le e^{(\lambda - s)(1 + \varepsilon)}$, and the theorem's own conclusion applied on $[0, s]$, $\|T_{y}(s) - T^{0}_{y}(s)\| \le s\,\varepsilon\,e^{s(1+\varepsilon)}$.

\emph{Derivative bounds.} The generators are $C^{1}$ in $y$: $\partial_{y}H_{y} = -\sum_{i}(\partial_{y}u^{i}_{y})p_{i}B_{i}$, so $\|\partial_{y}H_{y}\| \le |\partial_{y}u_{y}| \le c_{2}$ by Cauchy-Schwarz again, and $\partial_{y}\mathcal{E}_{y} = dE_{\tau}(\partial_{y}\gamma_{y})[u_{y}] + E_{\tau}(\gamma_{y})[\partial_{y}u_{y}]$, so $\|\partial_{y}\mathcal{E}_{y}\| = \delta$ at most by definition, and $\delta \le \varepsilon_{1}c_{1} + \varepsilon c_{2}$ when $\|dE_{\tau}\| \le \varepsilon_{1}$ and $\|\partial_{y}\gamma_{y}\| \le c_{1}$, using linearity of $E_{\tau}$ in the direction. For a coordinate family the first term is $\partial_{x}E_{\tau}[u]$ times $\partial_{y}x = 1$ and the second vanishes, so $\delta$ is the coordinate derivative of the error; in Hubble-normalised length a coordinate displacement has $c_{1}$ growing like a negative power of $\tau$, which is why the hypothesis is stated through $\delta$ and not through $\varepsilon_{1}$ and $c_{1}$ separately. By differentiable dependence of solutions of linear equations on parameters, $D := \partial_{y}T_{y}$ and $D^{0} := \partial_{y}T^{0}_{y}$ exist and solve $D' = GD + (\partial_{y}G)T$, $(D^{0})' = HD^{0} + (\partial_{y}H)T^{0}$, with $D(0) = D^{0}(0) = 0$. For the model, variation of constants gives $D^{0}(s) = \int_{0}^{s}T^{0}(s)T^{0}(\sigma)^{-1}(\partial_{y}H)(\sigma)T^{0}(\sigma)\,d\sigma$, and since $\|T^{0}(s)T^{0}(\sigma)^{-1}\| \le e^{s - \sigma}$ and $\|T^{0}(\sigma)\| \le e^{\sigma}$, $\|D^{0}(s)\| \le c_{2}\,s\,e^{s}$.

\emph{The difference.} $X := D - D^{0}$ satisfies
\[
X' \;=\; G X \;+\; \mathcal{E}\,D^{0} \;+\; (\partial_{y}\mathcal{E})\,T \;+\; (\partial_{y}H)\,(T - T^{0}), \qquad X(0) = 0,
\]
as one sees by subtracting the two derivative equations and writing $G D - H D^{0} = G X + (G - H)D^{0}$ and $(\partial_{y}G)T - (\partial_{y}H)T^{0} = (\partial_{y}\mathcal{E})T + (\partial_{y}H)(T - T^{0})$. Variation of constants gives $X(\lambda) = \int_{0}^{\lambda}T(\lambda)T(s)^{-1}\,[\cdots](s)\,ds$, and the bracket is bounded by $\varepsilon c_{2}s e^{s} + \delta\,e^{s(1+\varepsilon)} + c_{2}\,s\,\varepsilon\,e^{s(1+\varepsilon)} \le e^{s(1+\varepsilon)}\bigl(\delta + 2\varepsilon c_{2}s\bigr)$. Multiplying by $e^{(\lambda - s)(1+\varepsilon)}$ and integrating over $[0, \lambda]$ gives $\|X(\lambda)\| \le e^{\lambda(1+\varepsilon)}\bigl(\lambda\delta + \varepsilon c_{2}\lambda^{2}\bigr)$, the display. For $c_{2} = 0$ only the $\lambda\delta$ term survives.

\emph{Consequences.} The right side depends on the family only through $\lambda$ and $c_{2}$, and on the metric only through $\varepsilon(\tau)$ and $\delta(\tau)$; $\varepsilon(\tau) \to 0$ by admissibility. If $\delta(\tau) \to 0$ the right side tends to $0$; with Theorem \claimref{interior.transport-level-convergence-along-tangential-curves} this is convergence of $T_{y}(\lambda)$ to $T^{0}_{y}(\lambda)$ in $C^{1}$ in $y$, uniformly, so the limit is $C^{1}$ in $y$ with the model's derivative (uniform convergence of functions and of their derivatives). The model transport $T^{0}_{y}(\lambda)$ depends on $y$ only through the directions $u_{y}$ and the exponents of Theorem \claimref{interior.weighted-limit-theorem}(3), so its derivative is explicit. Nothing in admissibility controls $\delta$: the definition bounds $E_{\tau}$ and $\|E_{\tau}\|/\tau$ only, and an error of the form $\tau^{1/2}\varphi(x/\tau)$ with $\varphi$ bounded and $\varphi'(0) \ne 0$ meets every decay clause while its coordinate derivative at $x = 0$ is $\tau^{-1/2}\varphi'(0)$, unbounded (the sharpness example filed alongside, synthesis.admissible-decay-does-not-control-tangential-derivatives, takes $\tau^{1/2}\varphi(x/\tau^{2})$, with coordinate derivative $\tau^{-3/2}$ at $x = 0$). So the conclusion is conditional on the stated hypothesis.

\emph{The check.} The script runs the discrete (Euler-product) form of this argument in exact arithmetic on $4 \times 4$ real matrices of $\mathfrak{so}(1,3)$, at the Schwarzschild exponents and at the Kasner point $(-\tfrac27, \tfrac37, \tfrac67)$: a family of curves with directions on the unit circle in its rational parametrisation and a noncommuting polynomial error field, differentiated in $y$ exactly by dual numbers. It verifies the discrete variation-of-constants identity for $T - T^{0}$ in value and in $y$-derivative, the model bound $\|H\| \le 1$ on rational unit directions in the max-row-sum norm, the chain-rule bound $\delta \le \varepsilon_{1}c_{1} + \varepsilon c_{2}$ with the four constants computed separately, and the discrete forms of both bounds, with $(1 + h(1 + \varepsilon))^{n}$ in place of $e^{\lambda(1+\varepsilon)}$. Two tightness controls show each term is needed: an error with no base-point dependence at a point where the direction is stationary has $\delta = 0$ but a nonzero derivative of the transport difference, so the $\varepsilon c_{2}\lambda^{2}$ term cannot be dropped; and over a curve of length $12$ the bound without the growth factor fails while the stated one holds.
